\Gls{dm} is an invisible form of matter whose subatomic nature remains
unknown, motivating searches for physics beyond the \gls{sm}. The \gls{ldmx}
searches for sub-GeV dark matter by comparing the momentum of an incoming
electron with the visible final state after a fixed-target interaction.
Missing momentum could indicate the production of invisible dark matter
particles, making accurate reconstruction of the beam electrons and their
detector activity essential
\cite{cirelli2026-darkmatter,akesson2018-ldmx,ldmx2025-overview}.

In experimental particle physics, \emph{reconstruction} uses detector
measurements to infer physically meaningful properties of the underlying
particle event. At increased beam intensity, several electrons can enter
\gls{ldmx} within the same readout window and produce overlapping showers in
the \gls{ecal} subdetector. This obscures which detector activity originated
from each electron. Rejecting every such event would discard part of the
delivered beam, whereas reliable reconstruction could make a larger fraction
of recorded events useful and thereby improve the practical physics reach
\cite{ldmx2025-overview}.

This thesis investigates supervised \gls{ml} methods for separating
overlapping showers in simulated two-electron and three-electron events. The
principal hit-level task assigns each selected \gls{ecal} hit to the electron
responsible for the largest share of its deposited energy, as determined from
simulation truth. An event is represented as a variable-size collection of
\gls{ecal} hits and supplemented with reconstructed trigger-pad-track
observations from the \gls{ts} subdetector. Information from both subdetectors
can therefore be processed jointly through representations learned by the
models, without projecting the event onto a fixed image.

Two controlled baseline models are studied: a \gls{gnn} using GravNet
convolutions and a Transformer encoder. Both receive \gls{ecal} hits and
reconstructed trigger-pad-track observations, with the value of the latter
evaluated through inference-time ablation. Both architectures learn relations within unordered,
irregular detector data. Their use is motivated by recent machine-learned
particle-flow methods that interpret heterogeneous detector information with
graph-based or Transformer-based models
\cite{qasim2019-gravnet,pata2021-mlpf,pata2026-mlpf}. A secondary exploratory
multitask Transformer also predicts energy fractions and electron count
alongside hit assignment, following an approach inspired by particle-flow
reconstruction.

The study asks how reliably the models can assign hits in two-electron and
three-electron pile-up events, which shower properties limit the separation,
whether \gls{ts} inputs add useful context beyond the \gls{ecal},
and how GravNet and Transformer encoders compare in reconstruction quality and
computational properties.

The scope is limited to offline reconstruction of simulated inclusive
non-signal events. It does not perform a \gls{dm} search or deploy a real-time
trigger algorithm. This thesis presents the reusable
\texttt{ml\_ldmx} platform developed for the study, which prepares LDMX
simulation data, trains models and produces hit-level and event-level
diagnostics \cite{ml_ldmx-github}.

\clearpage
\subsection{Aim and Study Strategy}

The work was organised around the following objectives:

\begin{itemize}
    \item \textbf{Investigate shower separation:} determine how reliably
    supervised models can assign ECal hits to their originating electrons in
    simulated \(2e\) and \(3e\) pile-up events.

    \item \textbf{Build a reusable research platform:} develop the
    \texttt{ml\_ldmx} pipeline for extracting LDMX simulation data,
    constructing truth targets, training models and evaluating predictions.

    \item \textbf{Compare model architectures:} evaluate modern set-based neural networks choices GravNet and Transformer baselines using common inputs and targets, considering both
    reconstruction quality and practical computational properties.

    \item \textbf{Identify limiting conditions:} relate reconstruction
    performance to electron multiplicity, shower overlap, ECal depth and model
    confidence.

    \item \textbf{Assess heterogeneous detector information:} test whether
    reconstructed trigger-pad observations provide useful context beyond the
    ECal measurements.

    \item \textbf{Explore a broader reconstruction task:} extend the fixed electron-count hit assignment with a mixed-multiplicity Transformer inspired by a machine learning-based particle flow algorithm from CMS that predicts electron
    count, contributing-electron sets and energy fractions.
\end{itemize}

The intended outcome is an offline proof-of-concept and a documented basis for
continued reconstruction studies, rather than a production-grade trigger algorithm,
a complete detector reconstruction chain or a dark-matter sensitivity study.

\glsresetall